% ch09_conclusion.tex : Conclusion, Limitations, and Future Work
% ACO-SLAM: Swarm Intelligence-Based Search and Rescue Drone Navigation
% Language: Hebrew (XeLaTeX)

\selectlanguage{hebrew}

\section{סיכום, מגבלות ועבודה עתידית}
\label{sec:ch09_conclusion}

פרק זה מסכם את תרומות המאמר, דן במגבלות המערכת המוצעת, ומציע כיווני מחקר עתידיים להרחבת היכולות והתמודדות עם האתגרים שנותרו.

\subsection{סיכום התרומות}
\label{sec:ch09_contributions}

מסגרת \en{ACO-SLAM} שהוצגה במאמר זה מדגימה שיפורים משמעותיים לעומת שיטות קיימות בניווט אוטונומי של להקות רחפנים למשימות חיפוש והצלה. ארבע תרומות מרכזיות הוצגו ואומתו באמצעות סימולציות מקיפות וניסויי מעבדה.

התרומה הראשונה היא מסגרת \en{ACO-SLAM} חדשנית הממזגת מפות פרומונים עם גרידי תפוסה הסתברותיים. מיזוג זה מאפשר ניצול מיטבי של מידע היסטורי (פרומונים) לצד מידע עדכני (תפוסה) לצורך קבלת החלטות בזמן אמת. תוצאות הסימולציה הראו כי גישה זו משיגה כיסוי מרחבי של $94.8\%$, שיפור של $8.5\%$ לעומת \en{Swarm-SLAM} \cite{Lajoie2023SwarmSLAM}.

התרומה השנייה היא אופטימיזציית גרף תנוחות מבוזרת עם קונצנזוס בהשראת נמלים. אלגוריתם קונצנזוס הנמלים משתמש בנמלים וירטואליות הנושאות השערות טרנספורמציה בין רחפנים, וממצע השערות אלו משוקללות לפי עקביותן עם המפה הנוכחית. גישה זו משיגה \en{ATE} של $5.4$ ס"מ, שיפור של $34\%$ לעומת \en{Swarm-SLAM}, תוך שימוש ב-$28$ סבבי תקשורת בלבד לעומת $45$ ב-\en{Swarm-SLAM}.

התרומה השלישית היא מיזוג חיישנים רב-מודאלי בהשראה ביולוגית המשלב \en{LiDAR}, מצלמה ו-\en{IMU} באמצעות שקלול מבוסס פרומונים. כל חיישן מחזיק עקבה פרומון המשקפת את מהימנותו העדכנית, ומשקל המיזוג מסתגל בזמן אמת לתנאי החיישן. שיטה זו משיגה \en{RMSE} מיקום של $3.5$ ס"מ וזמן התאוששות מנפילת חיישן של $1.2$ שניות בלבד, לעומת $5.2$ ס"מ ו-$2.3$ שניות בשיטת \en{EKF} הסטנדרטית \cite{Kalman1960}.

התרומה הרביעית היא תכנון מסלול אדפטיבי בזמן אמת המותאם למשימות \en{SAR}. קצב התאיידות אדפטיבי $\rho(t)$ היורד ככל שמתגלים ניצולים, ומשקל חקר אדפטיבי $\alpha(t)$ הגדל עם הזמן, מאפשרים איזון דינמי בין חקר מרחבי לבין ניצול מידע על מיקומי ניצולים. אסטרטגיה זו משיגה $8.5$ ניצולים בממוצע וזמן משימה של $23.2$ דקות, שיפור של $30.8\%$ ו-$23.7\%$ בהתאמה לעומת \en{Swarm-SLAM}.

\subsection{מגבלות}
\label{sec:ch09_limitations}

למרות התוצאות המעודדות, למסגרת המוצעת מספר מגבלות משמעותיות הדורשות התייחסות במחקר עתידי.

המגבלה הראשונה היא יכולת ההרחבה ללהקות גדולות מ-50 רחפנים. עלות התקשורת הכוללת גדלה כ-$\mathcal{O}(N^{2})$ עבור העברת מפות מלאות, כאשר $N$ הוא מספר הרחפנים. עבור $N = 50$, רוחב הפס הנדרש להעברת מפות פרומונים מלאות הוא כ-10 מגה-ביט לשנייה, העולה על הקיבולת של רשת \en{WiFi} אד-הוק טיפוסית. אסטרטגיית רחפן-שליח ודחיסת מפות מפחיתות בעיה זו אך אינן פותרות אותה לחלוטין עבור להקות גדולות מאוד.

המגבלה השנייה היא אילוצי זמן אמת על מעבדי \en{ARM-class} הנפוצים ברחפנים קטנים. זמן העיבוד הממוצע למחזור \en{SLAM} (10 הרץ) היה $12.3$ אלפיות שנייה על גבי \en{NVIDIA Jetson Orin NX}, אך על מעבדים חלשים יותר (כגון \en{Raspberry Pi 4}) זמן העיבוד צפוי להיות ארוך משמעותית, העלול לגרום לעיכובים בלולאת הבקרה.

המגבלה השלישית היא התדרדרות חיישנים בתנאי עשן, אבק או ערפל האופייניים לסביבות אסון. מיזוג החיישנים המבוסס פרומונים מספק התדרדרות הדרגתית תחת אובדן חיישנים, אך במקרים של אובדן מוחלט של כל החיישנים האופטיים (\en{LiDAR} ומצלמה), המערכת מסתמכת על \en{IMU} בלבד ומתדרדרת במהירות.

המגבלה הרביעית היא ההנחה של סביבה סטטית יחסית. האלגוריתם מניח כי מפת התפוסה אינה משתנה במהלך המשימה, הנחה שאינה תמיד נכונה בסביבות אסון דינמיות (לדוגמה, מבנים מתמוטטים, מפולות).

\begin{equation}
\lim_{N \to \infty} \text{CommunicationOverhead}(N) = \mathcal{O}(N^{2})
\label{eq:ch09_scalability}
\end{equation}

\subsection{כיווני מחקר עתידיים}
\label{sec:ch09_future_work}

שלושה כיווני מחקר עתידיים עיקריים מזוהים להרחבת היכולות והתמודדות עם המגבלות שזוהו.

הכיוון הראשון הוא שימוש בלמידת חיזוק עמוקה (\en{Deep Reinforcement Learning}) לכיוונון אוטומטי של פרמטרי \en{ACO}. במקום לקבוע את הפרמטרים $\alpha$, $\beta$, $\rho$ באופן ידני או באמצעות סריקת רשת, סוכן \en{DRL} ילמד להתאים אותם באופן דינמי לתנאי הסביבה המשתנים. מחקרים ראשוניים בתחום זה \cite{Nedjah2022DRL} מראים כי גישת \en{DRL} יכולה לשפר את ביצועי החקר בעד $15\%$ נוספים. סוכן ה-\en{DRL} יקבל כקלט את מצב מפת הפרומונים, מיקומי הרחפנים, היסטוריית הגילויים, ויפיק כפלט את ערכי הפרמטרים האופטימליים למחזור הבא.

הכיוון השני הוא שימוש בלהקות רחפנים הטרוגניות עם חיישנים ייעודיים. להקה הטרוגנית תשלב רחפני \en{LiDAR} למיפוי מדויק, רחפני מצלמה רב-ספקטרלית לזיהוי ניצולים דרך עשן ואבק, רחפני סונאר לזיהוי ניצולים מתחת להריסות, ורחפני ממסר לשמירת קישוריות תקשורת. גישה זו תאפשר התמודדות טובה יותר עם תנאי סביבה קשים ותגדיל את סיכויי הזיהוי של ניצולים.

הכיוון השלישי הוא הרחבת מפות הפרומונים למרחב תלת-ממדי מלא עבור \en{SAR} רב-קומתי. מפת פרומונים תלת-ממדית $\tau_{ijk}(t)$ מאפשרת ייצוג של מידע חקר בבניינים רבי-קומות, מנהרות תת-קרקעיות, וסביבות תלת-ממדיות מורכבות אחרות. כלל עדכון הפרומון התלת-ממדי הוא:

\begin{equation}
\tau_{ijk}(t+1) = (1-\rho)\tau_{ijk}(t) + \sum_{d \in \mathcal{D}} \Delta\tau_{ijk}^{(d)}(t)
\label{eq:ch09_3d_pheromone}
\end{equation}

כאשר $\tau_{ijk}(t)$ הוא ערך הפרומון בתא תלת-ממדי $(i,j,k)$ ו-$\mathcal{D}$ היא קבוצת הרחפנים. הרחבה זו תאפשר למערכת לפעול בסביבות מורכבות יותר, כגון בניינים רבי-קומות לאחר רעידת אדמה.

\begin{table}[htbp]
\centering
\caption{סיכום מגבלות המערכת המוצעת ופתרונות עתידיים מוצעים.}
\label{tab:ch09_limitations}
\adjustbox{max width=\columnwidth}{%
\begin{tabular}{lcc}
\toprule
\textbf{מגבלה} & \textbf{חומרה} & \textbf{פתרון עתידי} \\
\midrule
הרחבה ל-$>50$ רחפנים & גבוהה & קיבוץ היררכי \\
זמן אמת על \en{ARM} & בינונית & האצת \en{FPGA} \\
התדרדרות חיישנים & גבוהה & יתירות רב-מודאלית \\
סביבות תלת-ממד & בינונית & מפת פרומון תלת-ממדית \\
סביבות דינמיות & בינונית & \en{SLAM} דינמי \\
\bottomrule
\end{tabular}%
}
\end{table}

\subsection{סיכום}
\label{sec:ch09_summary}

לסיכום, \en{ACO-SLAM} מהווה צעד משמעותי קדימה בניווט אוטונומי של להקות רחפנים למשימות חיפוש והצלה, תוך שילוב עקרונות ביולוגיים עם כלים הסתברותיים מתקדמים. המערכת מדגימה שיפורים של $34\%$ ב-\en{ATE}, $8.5\%$ בכיסוי מרחבי, $30.8\%$ בניצולים שנמצאו, ו-$23.7\%$ בזמן משימה לעומת שיטות קיימות. מיזוג החיישנים בהשראה ביולוגית מספק התדרדרות הדרגתית תחת אובדן חיישנים, ואלגוריתם קונצנזוס הנמלים משיג התכנסות מהירה יותר מאופטימיזציה מבוזרת סטנדרטית.

המגבלות העיקריות כוללות יכולת הרחבה ללהקות גדולות, אילוצי זמן אמת על חומרה מוגבלת, והתדרדרות חיישנים בתנאי סביבה קשים. כיווני המחקר העתידיים המוצעים, ובהם למידת חיזוק עמוקה לכיוונון פרמטרים, להקות הטרוגניות, ומפות פרומונים תלת-ממדיות, צפויים לאפשר פריסה מבצעית של להקות רחפנים אוטונומיות בסביבות מורכבות ומשתנות.

הפיתוח העתידי של מערכת \en{ACO-SLAM} צפוי לתרום להצלת חיים במצבי אסון, על ידי אפשור חיפוש מהיר, יעיל ובטוח יותר של ניצולים בסביבות מסוכנות ובלתי נגישות לבני אדם.

\cite{Coppola2020Swarming, Chung2018Swarm, Nedjah2022DRL, Lajoie2023SwarmSLAM, Kalman1960, Yan2020Survey}

\appendix
\section{רשימת סמלים ומשתנים}
\label{sec:ch09_appendix}

\begin{table*}[htbp]
\centering
\caption{רשימת סמלים ומשתנים במערכת \en{ACO-SLAM}.}
\label{tab:ch09_symbols}
\adjustbox{max width=\columnwidth}{%
\begin{tabular}{lll}
\toprule
\textbf{סמל} & \textbf{הגדרה} & \textbf{יחידות} \\
\midrule
$\mathbf{x}_{k}^{(i)}$ & וקטור מצב של רחפן $i$ בזמן $k$ & משתנה \\
$\mathbf{p}_{k}^{(i)}$ & מיקום רחפן $i$ בזמן $k$ & מטר \\
$\mathbf{q}_{k}^{(i)}$ & כיוון רחפן $i$ בזמן $k$ (קווטרניון) & --- \\
$\mathbf{v}_{k}^{(i)}$ & מהירות רחפן $i$ בזמן $k$ & מטר/שנייה \\
$\mathbf{u}_{k}^{(i)}$ & וקטור בקרה של רחפן $i$ בזמן $k$ & משתנה \\
$\mathbf{z}_{k}^{(i)}$ & מדידת חיישן של רחפן $i$ בזמן $k$ & משתנה \\
$\mathcal{M}_{t}$ & מפה גלובלית בזמן $t$ & --- \\
$\tau_{ij}(t)$ & ערך פרומון בתא $(i,j)$ בזמן $t$ & --- \\
$\rho$ & קצב התאיידות פרומון & --- \\
$\alpha$ & משקל פרומון בהסתברות מעבר & --- \\
$\beta$ & משקל היוריסטי בהסתברות מעבר & --- \\
$\omega$ & משקל תגמול ניצול & --- \\
$\eta_{ij}(t)$ & מידע היוריסטי לתא $(i,j)$ בזמן $t$ & --- \\
$P_{ij}^{(k)}(t)$ & הסתברות מעבר של רחפן $k$ מתא $i$ לתא $j$ & --- \\
$\mathcal{G}(t)$ & גרף תקשורת בזמן $t$ & --- \\
$d_{\text{comm}}$ & טווח תקשורת מרבי & מטר \\
$\mathbf{T}_{ij}$ & טרנספורמציה יחסית בין רחפן $i$ לרחפן $j$ & מטר, רדיאן \\
$\text{ATE}$ & שגיאת מסלול מוחלטת & מטר \\
$\text{RPE}$ & שגיאת תנוחה יחסית & מטר/מטר \\
$P_{\text{victim}}(\mathbf{p}, t)$ & הסתברות לניצול במיקום $\mathbf{p}$ בזמן $t$ & --- \\
$\mathcal{V}_{ij}(t)$ & תגמול ניצול בתא $(i,j)$ בזמן $t$ & --- \\
$\mathcal{I}_{ij}(t)$ & תוספת מידע בתא $(i,j)$ בזמן $t$ & --- \\
$w_{s}^{(i)}(t)$ & משקל מיזוג חיישן $s$ על רחפן $i$ בזמן $t$ & --- \\
$\mathbf{K}_{k}$ & רווח \en{Kalman} בזמן $k$ & --- \\
$\mathbf{\Sigma}_{ij}$ & מטריצת שיתוף פעולה של אילוץ $(i,j)$ & --- \\
$\mathcal{R}(t)$ & קבוצת רחפני ממסר בזמן $t$ & --- \\
$\Delta_{\tau}$ & גודל צעד כימות פרומון & --- \\
\bottomrule
\end{tabular}%
}
\end{table*}

\cite{Coppola2020Swarming, Chung2018Swarm, Nedjah2022DRL, Lajoie2023SwarmSLAM, Kalman1960, Yan2020Survey}